\chapter{Single Fields}

\section{Electrostatics}
\subsection{Results}
The electric field intensity vector $\bE$ can be derived from the
nodal degree of freedom, the electric potential $V$ by applying a
gradient operator. This element quantity is calculated by applying
the config-file-commands

\begin{verbatim}
<storeResults>
   <elemResults type="elecFieldIntensity" region="elecst2d"/>
</storeResults>
\end{verbatim}

In the command {\bf $<$region$>$}, all areas in which the electric field has to
be calculated, must be listed. If no region  is specified, the elecric
field will be calculated in all regions where the PDE is defined. To
calculate the electric energy, one has to use the command 

\begin{verbatim}      
<storeResults>
   <elemResults type="elecEnergy" region="reg2"/>
</storeResults>
\end{verbatim}
in the config-file. The result is written to the info-file.


\subsection{Example}

{\footnotesize
\begin{verbatim}
<?xml version="1.0"?>

<cfsSimulation xmlns="http://www.cfs++.org">

  <analysis type="static"/>
  <geometry type="axi"/>
  <input format="mesh"/>
  <output format="unv"/>
  <materialData file="mat.dat"/>

  <domain>
    <region name="gndEl"  material="steel"/>
    <region name="hotEl"  material="steel"/>
    <region name="dielEl" material="silicon"/>
    <region name="airEl"  material="air"/>

    <nodes name="hot"/>
    <nodes name="gnd"/>
  </domain>

  <pdeList>
    <electrostatic>
       <region name="gndEl"/>
       <region name="hotEl"/>
       <region name="dielEl"/>
       <region name="airEl"/>

      <bcsAndLoads>
        <dirichletHom name="gnd"/>
        <dirichletInhom name="hot" value="1.0"/>
      </bcsAndLoads>

      <storeResults>
        <elemResults type="elecFieldIntensity"/>
        <elemResults type="elecEnergy"/>
      </storeResults>

    </electrostatic>
  </pdeList>

</cfsSimulation>
\end{verbatim}
}


\newpage
\section{Mechanics}
\subsection{Degree of freedoms}
In the 3D case we have {\bf ux, uy, uz} and in the 2D case
(axisymmetric as well as plane) {\bf ux, uy}.


\subsection{Subtypes}

A specification of the geometric type of the problem is done on the top-level
via the \xml{geometry} element, which allows the values \xml{3d}, \xml{plane}
and \xml{axi} for its \xml{type} attribute. For a mechanic PDE this definition
must be repeated in the \xml{mechanic} element via its \xml{subtype} entry.
\begin{verbatim}
       <mechanic subType="stype">
\end{verbatim}
\noindent Here stype can currently take one of the following three values

\begin{itemize}
\item planeStrain \emph{(the default)}
\item 3d
\item axi
\end{itemize}

\emph{Note:} The \xml{subType} value must of course match to the \xml{type}
value of the geometry specification and a mixture of different subTypes within
one analysis is thus not possible.

\subsection{Node to Ground Springs}
In mechanic it is possible to consider node to ground springs. In a
static analysis one can consider a stiffness value and
in a transient analysis additionally a mass and a damping value.
With the parameter name an existing selection of nodes must be given
as it is the case for loads.
\begin{verbatim}      
<bcsAndLoads>
 <spring name="springNodes" dof="uy" massValue="1" dampingValue="1" 
    stiffnessValue="1"/>
</bcsAndLoads>
\end{verbatim}

\subsection{Nodal Results}
Possible nodal results are

\begin{verbatim}
   <nodeResults type="mechDisplacement" region="bar"/>
   <nodeResults type="mechVelocity" region="bar"/>
   <nodeResults type="mechAcceleration" region="bar"/>
\end{verbatim}

\subsection{Element Results}
The command for calculating the mechanical stresses is
\begin{verbatim}      
<storeResults>
   <elemResults type="mechStress" region="bar"/>
</storeResults>
\end{verbatim}

\subsection{Example}

{\footnotesize 
\begin{verbatim}
<?xml version="1.0"?>
<cfsSimulation xmlns="http://www.cfs++.org">

  <analysis type="static"/>
  <geometry type="plane"/>
  <input format="mesh"/>
  <output format="unv"/>
  <materialData file="mat.dat"/>

  <domain>
    <region name="cantilev"  material="steel"/>
    <nodes name="leftN"/>
    <nodes name="rightN"/>
  </domain>

  <pdeList>
    <mechanic subtype="planeStrain">
       <region name="cantilev"/>
       <bcsAndLoads>
         <dirichletHom name="leftN" dof="ux"/>
         <dirichletHom name="leftN" dof="uy"/>
         <dirichletHom name="rightN" dof="ux"/>
         <load         name="rightN" dof="uy" value="-1.0"/>
       </bcsAndLoads>
    </mechanic>
  </pdeList>

</cfsSimulation>
\end{verbatim}
}


\newpage
\section{Acoustics}

\subsection{General Setup}
The \xml{acoustic} element is appended by the following child elements
\begin{itemize}
\item \xml{region}
\item \xml{surface}
\item \xml{timeSteppingFormulation} effective stiffness formulation 
is default(optional), effective mass matrix formulation can be chosen unless one is not using the fractional damping model
\item \xml{bcsAndLoads} (optional)
\item \xml{nonLinear}\footnote{Other Nonlinearities than nonlinear acoustic
 PDE, which can be chosen by an attribute of \xml{region} element. This is
 currently not implemented.}
\item \xml{storeResults} (optional)
\end{itemize}
and the following attributes
\begin{itemize}
\item \xml{formulation} used to choose between an acoustic pressure or
      acoustic potential PDE formulation, specified by either choosing
      either \xml{acouPressure} or \xml{acouPotential} as value.
\item \xml{subType} for flownoise computations set this to \xml{flowNoise}.
\end{itemize}

\subsection{Formulation}
For linear acoustics, the PDE in pressure and acoustic potential has the same 
form. Allthough default is acoustic potential. If acoustics is coupled to 
other PDEs, acoustic potential is obligatory. Westervelt's nonlinear PDE is 
in pressure, Kuznetsov's nonlinear PDE is in acoustic potential.

\subsection{Damping and nonlinear acoustic simulation}
Damping and nonlinear acoustic PDE formulations can be chosen for each 
region individually. Damping is a child element of \xml{region},
type is either no damping, rayleigh damping (which is not physical), 
thermoviscous damping and fractional damping according to a 
frequency power law.
\xml{nonLinear} comes as attribute to region. It can take the value \xml{no},
\xml{westervelt}, or \xml{kuznetsov}.

\subsection{Boundary Conditions}
Unique to the acoustic PDE is the possibiliy to prescribe inhomogeneous Neumann boundary conditions. In acoustic potential formulation this means a normal velocity on a surface. Remember, that the specified surface elements have to appear in the \xml{domain} section and in the \xml{acoustic} section below the region definitions. The syntax is:

\begin{verbatim} 
<bcsAndLoads>
  <neumannInhom name="transducer" value="1e-6" dynamics="sweep5M.dat"/>
</bcsAndLoads>
\end{verbatim}

\subsection{Nodal Results}
The config-file command for storing e.g. the acoustic potential in the result
file is

\begin{verbatim}      
<storeResults>
   <nodeResults type="acouPotential" region="fluid1"/>
</storeResults>
\end{verbatim}

Analogously one can specify storing of certain node values with the help of the \xml{nodeHistory} element.

\begin{verbatim}    
<storeResults>
  <nodeHistory saveNodes="historyPoints" type="acouPotential"/>
</storeResults>
\end{verbatim}

Remember, that the quantity, you want to store as result, depends on the formulation of the PDE. It is straightforward to retrieve the quantity acoustic potential, when the formulation of the PDE is in acoustic potential. But if your formulation is in acoustic pressure, it is nonsense to retrieve the quantity acoustic potential, because this would need integration. A combination, which does work, is formulation in acoustic potential and results in acoustic pressure. This will be discussed in the next section.

\subsection{Element Result}

Because $p = \rho \partial \psi / \partial t$, and we can only assign a certain density to an element in the computational domain, pressure is an element result for acoustic potential formulation. It is specified with

\begin{verbatim}
<storeResults>
  <elemResults type="acouPressure" />
  <elemHistory type="acouPressure" saveElems="receiver"/>
</storeResults>
\end{verbatim}

%\subsubsection{Element Velocity $\{v\}$}
%\begin{print}
%  The velocity can be calculated as an element result by applying an
%  {\em attgroup} command.
%\end{print}
%\begin{verbatim}
%attgroup,1,calc,calc_gvp
%\end{verbatim}


\subsection{Absorbing Boundary Conditions (ABC)} 
In the config file, the command

\begin{verbatim}    
<bcsAndLoads>  
   <absorbingBCs name="abc"/>
</bcsAndLoads>
\end{verbatim}
      
is needed, which defines the boundaries at which the ABCs are working.
It should be clear, that the dimension of the boundary elements are
one order lower compared to the dimension of the "domain" elements.
Therefore, the boundary elements have to be listed with the config-file
command

\begin{verbatim}
<domain>
   <surface name="abc"/>
</domain>
\end{verbatim}

\subsection{Aeroacoustics}
In aeroacoustics, sound sources are computed from a CFD simulation. This is
done by another tool, called \texttt{cplreader} (see the \texttt{cplreader}
manual for details). The sound sources are then used as a load on the right
hand side of the acoustic PDE. Therefore the corresponding XML element is
called \texttt{rhsValues}. Here is an example:

\begin{verbatim}
<rhsValues region="fluid" inputId="acouGrid">
  <factor>1.0</factor>
  <asyncSteps>nearest</asyncSteps>
  <interpolation justInterpolate="yes">
    <srcRegions names="partition1" inputId="sourceGrid" coordSysId="default"/>
    <xyPlane z="0.0" tol="1.0e-10"/>
    <tolerances>
      <global start="1e-4" end="1e-3" inc="9e-4"/>
      <local start="1e-3" end="5.1e-2" inc="1e-2"/>
    </tolerances>
    <restartFileMode>rw</restartFileMode>
    <nodeWarnings display="verbose" writeNodes="yes"/>
  </interpolation>
</rhsValues>
\end{verbatim}

The attribute \texttt{region} identifies the region on which the source term
is to be applied; it is always required. If there is no sub-element
\texttt{interpolation} then source terms are read from the file given through
the attribute \texttt{inputId} (which defaults to \texttt{default}).

Inside the \texttt{factor} element one can supply a mathematical expression
(see chapter~\ref{chap:MathExpr}) which is multiplied with the source terms.
The expression may depend on time or frequency, and space. This is especially
useful in a transient simulation that begins after the flow field has had some
time to develop turbulence. In this case, immediately switching on the full
source term can lead to numerical errors. A ramp function has proven to give
much better results. A commonly used ramp is
\verb|if( t lt T; sin(2*pi*t/(4*T))^2; 1)|, where \texttt{T} must be replaced
with the time duration of the ramp.

The \texttt{asyncSteps} element tackles that fact that time steps in CFD are
usually smaller than in computational acoustics. Setting this value to
\texttt{nearest} selects source terms from the time step that is nearest to the
one that is currently being computed. \texttt{interpolate} will result in a
linear temporal interpolation of source terms between the two time steps that
the current time lies in between. In order to use the default behavior of
CFS++, selecting data by step number, use \texttt{no} here.

The \texttt{interpolation} element tells CFS++ to interpolate source terms from
a different grid. This is required, if a CFD grid is not suitable for
acoustics. The attribute \texttt{justInterpolate} may be set to \texttt{yes} if
just the interpolation should be computed, but not the acoustic field. The
\texttt{srcRegions} element defines the CFD grid; \texttt{names} is the list
of regions (separated by space, \verb|;| or \verb#|#) that contain
source terms; \texttt{inputId} identifies the input file from which to read the
CFD grid and source terms; optionally \texttt{coordSysId} may specify the
coordinate system of the CFD grid if it is different from the acoustics grid.

The \texttt{xyPlane} element is only relevant if the acoustics grid is
two-dimensional, but the CFD grid is three-dimensional. Then one may specify
the \texttt{z} coordinate of the x-y-plane (defaults to 0). Additionally
\texttt{tol} gives the tolerance of \texttt{z} in meters, so that nodes in the
CFD grid need not lie exactly on the plane.

Inside the \texttt{tolerances} element one may specify a range for the global
and local tolerances, which are used to determine the interpolation weights.
The global tolerance (in meters) determines the space around a node in the CFD
grid, when its surrounding element in the acoustics grid is searched for. The
local tolerance is the space around an element (in the element's local
coordinate system) where nodes are still accepted, although they lie outside of
the element. These tolerance are most relevant for domains with curved
boundaries, which are discretized very differently by CFD grid and acoustics
grid. Tolerance ranges are swept in order to use small tolerances wherever
possible and large tolerances only in places where it is necessary.

\texttt{restartFileMode} determines if a restart file containing the
interpolation weights is to be read (\texttt{r}), written (\texttt{w}) or both
(\texttt{rw}). In this way the interpolation weights need to be computed only
once and can be reused if a simulation is restarted.

The \texttt{nodeWarnings} element defines how nodes are handled which could not
be mapped during interpolation. Any message is suppressed if the attribute
\texttt{display} is set to \texttt{no}. The total number of missed nodes is
reported, when set to \texttt{yes}. A message for each single node is generated
by a value of \texttt{verbose}. Optionally a text file containing the node
numbers is written, if the attribute \texttt{writeNodes} is set to \texttt{yes}.

\subsection{Example}

{\footnotesize 
\begin{verbatim}
<?xml version="1.0"?>
<cfsSimulation xmlns="http://www.cfs++.org">
  <analysis type="transient"/>
  <geometry type="plane"/>
  <domain>
    <region material="softTissue" name="reg1"/>
    <region material="bone" name="reg2"/>
    <nodes name="vp-restr"/>
  </domain>
  <pdeList>
    <acoustic formulation="acouPotential">
      <region name="reg1" nonLinear="no">
        <damping type="no"/>
      </region>
      <region name="reg2" nonLinear="no"/>
      <bcsAndLoads>
        <dirichletInhom dynamics="sin1M40cos.dat" name="vp-restr" value="1.0"/>
      </bcsAndLoads>
      <storeResults>
        <nodeHistory saveNodes="historyPoints" type="acouPotential"/>
      </storeResults>
    </acoustic>
  </pdeList>
  <transient>
    <numSteps>    600     </numSteps>
    <firstDt>     2.5e-08 </firstDt>
    <stepSaveBeg>   1     </stepSaveBeg>
    <stepSaveEnd> 600     </stepSaveEnd>
    <stepSaveInc>   1     </stepSaveInc>
    <timeDataFile name="sin1M40cos.dat"/>
  </transient>
</cfsSimulation>
\end{verbatim}
}


\newpage
\section{Magnetics}
\subsubsection{Element Results}
In the case of magnetic field calculation, many element results can be
deduced by a postprocessing step:

\begin{itemize}
\item Magnetic induction (magnetic flux density) $\bB$
    \newline Derived from the magnetic vector potential $\bA$ by
    applying the curl operator.
  
\begin{verbatim}
<storeResults>
   <elemResults type="magFluxDensity" region="yoke"/>
</storeResults>
\end{verbatim}
  
\item Eddy current density $\bJ_{\rm e}$ 
    \newline Currents induced by a magnetic field in a conductive
    material. 

\begin{verbatim}
<storeResults>
   <elemResults type="magEddyCurrent" region="yoke"/>
</storeResults>
\end{verbatim}
  
  
\item Magnetic energy $W$ 
    \newline Measure for the energy stored in the magnetic assembly.
  
\begin{verbatim}
<storeResults>
   <elemResults type="magEnergy" region="yoke"/>
</storeResults>
\end{verbatim}

\item  Inductivity $L$ 

\begin{verbatim}
<saveFileL>   lFile.dat   </saveFileL>
\end{verbatim}

The inductivity $L$ of the magnetic assembly, as well as
the magnetic flux $\Phi$ are calculated and stored in a file with
the name "lFile.dat"  given in column oriented form:

{\em time inductivity magFlux}.  
  

%\item Magnetic force $F_{\rm V}$, {\em calc\_fv}, $<$emag$>$-file
%  \newline Used for calculating the magnetic force $F_{\rm V}$, due to
%  a current loaded coil within a magnetic field (Lorentz force)..

\end{itemize}

\subsection{Coil Definition}

The list of coils has to be defined inside the tags

\begin{verbatim}
<coils>

</coils>
\end{verbatim}

\subsubsection{Coil Types}

In our simulation, we distinguish different coil types:
\begin{itemize}
\item Measurement coils
  \begin{itemize}
  \item 2D: \verb! <measurementCoil2d name="coil1">!
  \item 3D: \verb! <measurementCoil3d name="coil2">!
  \end{itemize}
\item Volage driven coils 
  \begin{itemize}
  \item 2D: \verb! <voltageCoil2d name="coil3">!
  \item 3D: \verb! <voltageCoil3d name="coil4">!
  \end{itemize}
\item Current driven coils 
  \begin{itemize}
  \item 2D: \verb! <currentCoil2d name="coil5">!
  \item 3D: \verb! <currentCoil3d name="coil6">!
  \end{itemize}
\end{itemize}

There are several parameters which must be defined for a coil. Every
coil needs an attribute {\bf name="xyz"}, standing nearby the
definition of the coil:

\verb! <voltageCoil2d name="coil3">!

\noindent
Further on, the cross section of a winding has to be
specified for every coil:

\verb!<windingCrossSection> 1e-6 </windingCrossSection>!

\noindent
The number of windings $N$ results from the cross section of the wire. It
is calculated by the overall cross section of the coil $A_c$ devided
by the cross section of a winding $A_w$

\begin{equation}
  \label{eq:N}
  N=\frac{A_c}{A_w}
\end{equation}

\noindent
Depending on the coil type, several additional tags are needed. They
are described in the next sections.


\subsubsection{Measurement Coils}

\begin{itemize}
\item Coil voltage $U$ respectively coil current $I$

\begin{verbatim}
<saveFileU>  uiFile.dat  </saveFileU>
\end{verbatim}

  With this tag the voltage respectively the current of a
  special coil element can be calculated. It's stored in a file
  named {\em resultFile} in the column oriented form:

  {\em time voltage}.

\end{itemize}

\subsubsection{Current- and Voltage Driven Coils}
In this coil type, the amplitude, phase (in case of harmonic analysis)
and the input function (in case of dynamic analyses) of the driving
current respectively voltage has to be given:

\begin{itemize}
\item Transient case
\begin{verbatim}
<value>       10           </value>
<dynamics>    timefile.dat </dynamics>
\end{verbatim}

\item Harmonic case
\begin{verbatim}
<value>       10           </value>
<phase>       90           </phase>
\end{verbatim}
\end{itemize}

\noindent
The name of the time function as well as the sepcification of the
phase is optional. If it's not given, a constant current with an
amplitude of {\bf value} is assumed.  
Voltage driven coils additionally require the total ohmic resistance
of the coil for solving the network equation:

\verb!<resistance>  10 </resistance>!

\subsubsection{2D Coils}
For all types of 2D-coils, an identification number is
necessary:

\verb! <id>  2  <id>!

\noindent
Due to the 2D- or axisymmetric modeling, only the cross-section of the
coil is modeled. The coil-id is
used to identify coil cross-sections belonging together (fed by the
same current). 
The sign of the id also specifies the direction of the current.

\subsubsection{3D Current Driven Coils}
For 3D current driven coils, the direction of the current has to be
given explicitly. In case of a straight current flow, three
possibilities can be given:

\verb!<currentFlow> direction </currentFlow>!

\noindent
with {\bf direction} standing for one of 
\begin{itemize}
\item xDir
\item yDir
\item zDir
\end{itemize}
   
For rotational current flow (e.g. in 3D-modeled axisymmetric coils),
additionally the midpoint and the direction of the rotational axis
must be defined:

\begin{verbatim}
<rotational>
   <midPointX> 0 </midPointX>
   <midPointY> 0 </midPointY>
   <midPointZ> 0 </midPointZ>
   <rotAxisX>   1 </rotAxisX>
   <rotAxisY>   0 </rotAxisY>
   <rotAxisZ>   0 </rotAxisZ>
</rotational>
\end{verbatim}

In this example, the midpoint lies at the origin, and the coil
rotation axis is equal to the x-axis.

\subsubsection{Coil Definition Examples}

The full description of a voltage driven coil in 2D with a winding
cross section of $10^{-3}\, \rm m^2$, a dynamic excitation of 10\,V with
the dynamic behaviour saved in the file "time.dat", 
and 10\,$\Omega$ resistance could be given as
\begin{verbatim}
<voltageCoil2d name="coil1">
   <windingCrossSection> 0.001  </windingCrossSection>
   <value>                  10  </value>
   <dynamics>         time.dat  </dynamics>
   <resistance>             10  </resistance>
   <id>                      2  </id>
</voltageCoil2d>
\end{verbatim}

A 3D rotational coil with the rotational axis parallel to the y-axis
including the point (2,2,0) should be defined as

\begin{verbatim}
<currentCoil3d name="coil2">
  <windingCrossSection> 1e-5 </windingCrossSection>
  <value>                100 </value>
  <dynamics>        time.dat </dynamics>
  <rotational>
    <midPointX>  2  </midPointX>
    <midPointY>  2  </midPointY>
    <midPointZ>  0  </midPointZ>
    <rotAxisX>    0  </rotAxisX>
    <rotAxisY>    1  </rotAxisY>
    <rotAxisZ>    0  </rotAxisZ>
  </rotational>
</currentCoil3d>
\end{verbatim}


\subsection{Permanent Magnets}
The use of permanent magnets in the simulation is very simple. In the
magnetic section, the command

\begin{verbatim}
<magnets>
   <magnet name="perm" orientX="0" orientY="0" orientZ="1" />
</magnets>
\end{verbatim}
has to list all domains, which hold a permanent magnet. The
orientation and strength of the permanent magnetic field is given by
the magnetization vector $\bM$ giben in Tesla described by the three
values {\bf orientX}, {\bf orientY} and {\bf orientZ}

\begin{equation}
  \label{eq:M}
  \bM = \left(
    \begin{tabular}{r}
      orientX\\
      orientY\\
      orientZ
    \end{tabular}
\right) \; .
\end{equation}

\subsubsection{Example}
{\footnotesize
\begin{verbatim}
<?xml version="1.0"?>
<cfsSimulation xmlns="http://www.cfs++.org">

   <analysis type="static"/>
   <geometry type="plane"/>
   <materialData file="mat.dat"/>

   <domain>
      <region name="myCoil" material="air"/>
      <region name="iron"   material="iron"/>
      <region name="air"    material="air"/>
      <nodes name="bound"/>
   </domain>

   <pdeList>
      <magnetic>

         <region name="myCoil" />
         <region name="iron"   />
         <region name="air"    />

         <bcsAndLoads>
            <dirichletHom   name="bound"/>
         </bcsAndLoads>

         <coils>
            <currentCoil2d name="myCoil">
               <!-- whole coil area: 100*5 mm^2 with 20 windings -->
               <windingCrossSection>50e-6 </windingCrossSection>
               <value>       1            </value>
               <dynamics>    time.dat     </dynamics>
               <id>          1            </id>
            </currentCoil2d>
         </coils>

         <storeResults>
            <elemResults type="magFluxDensity" region="myCoil"/>
            <elemResults type="magFluxDensity" region="iron"/>
            <elemResults type="magFluxDensity" region="air"/>
         </storeResults>

      </magnetic>
   </pdeList>
</cfsSimulation>
\end{verbatim}
}

\newpage
\section{Heat Conduction}
\subsection{Results}
Until now, only the temperature, which is a nodal result, can be chosen for output either for the whole domain or at named history nodes.

\begin{verbatim}
<storeResults>
   <nodeResults type="heatTemperature"/>
</storeResults>
\end{verbatim}


\subsection{Example}

{\footnotesize
\begin{verbatim}
<?xml version="1.0"?>
<cfsSimulation xmlns="http://www.cfs++.org">

  <analysis type="transient"/>
  <geometry type="plane"/>
  <input format="mesh"/>
  <output format="gmv"/>
  <gmv>
    <fixedGrid> yes </fixedGrid>
    <binaryFormat> no </binaryFormat>
  </gmv>
  <materialData file="mat.dat"/>

  <domain>
    <region name="slab" material="steel" />
  </domain>

  <pdeList>
    <heatConduction>
      <region name="slab"/>

      <bcsAndLoads>
        <dirichletInhom name="bcleft" value="0.0"/>
        <dirichletInhom name="bcright" value="100.0" dynamics="sin80no.dat"/>
      </bcsAndLoads>

      <storeResults>
        <nodeResults type="heatTemperature"/>
        <nodeHistory type="heatTemperature" saveNodes="historyNode"/>
      </storeResults>
    </heatConduction>
  </pdeList> 

  <transient>
    <numSteps>     400 </numSteps>
    <firstDt>        2 </firstDt>
    <stepSaveBeg>    1 </stepSaveBeg>
    <stepSaveEnd>  400 </stepSaveEnd>
    <stepSaveInc>    1 </stepSaveInc>
    <timeDataFile name="sin80no.dat"/>
  </transient>

</cfsSimulation>
\end{verbatim}
}
